\begin{remark*}[Exposition]
Interval arithmetic replaces exact real inputs by intervals of possible values.
Given two intervals $A$ and $B$, each arithmetic operation collects all
pointwise results obtained by choosing one element from $A$ and one element
from $B$. The set-based definition guarantees enclosure of all possible
outcomes, while the endpoint formulas provide an operational way to compute
the resulting interval.
\end{remark*}

\begin{definition}[Set-based interval operation]
\cite{UnknownIntroductionIntervalAnalysis2002}
\label{def:set-based-interval-operation}
Let $\circ$ denote one of the four arithmetic operators $+$, $-$, $\cdot$, or $/$.
For $A,B\in \mathcal{I}(\mathbb{R})$, define
\[
A\circ B
=
\{\,a\circ b\mid a\in A,\ b\in B\,\},
\]
with $A/B$ undefined if $0\in B$.
\end{definition}

\begin{definition}[Interval addition]
\cite{UnknownIntroductionIntervalAnalysis2002}
\label{def:interval-addition}
For $A=[\underline a,\overline a]$ and
$B=[\underline b,\overline b]$ in $\mathcal{I}(\mathbb{R})$, define
\[
A+B
=
\{\,a+b\mid a\in A,\ b\in B\,\}
=
[\,\underline a+\underline b,\ \overline a+\overline b\,].
\]
\end{definition}

\begin{example}
If $A=[1,3]$ and $B=[2,5]$, then
\[
A+B=[1+2,3+5]=[3,8].
\]
\end{example}

\begin{definition}[Interval subtraction]
\cite{UnknownIntroductionIntervalAnalysis2002}
\label{def:interval-subtraction}
For $A=[\underline a,\overline a]$ and
$B=[\underline b,\overline b]$ in $\mathcal{I}(\mathbb{R})$, define
\[
A-B
=
\{\,a-b\mid a\in A,\ b\in B\,\}
=
[\,\underline a-\overline b,\ \overline a-\underline b\,].
\]
\end{definition}

\begin{example}
If $A=[1,3]$ and $B=[2,5]$, then
\[
A-B=[1-5,3-2]=[-4,1].
\]
\end{example}

\begin{definition}[Interval multiplication]
\cite{UnknownIntroductionIntervalAnalysis2002}
\label{def:interval-multiplication}
For $A=[\underline a,\overline a]$ and
$B=[\underline b,\overline b]$ in $\mathcal{I}(\mathbb{R})$, define
\[
A\cdot B
=
\{\,ab\mid a\in A,\ b\in B\,\}.
\]
Equivalently,
\[
A\cdot B
=
\Bigl[
\min\{
\underline a\,\underline b,
\underline a\,\overline b,
\overline a\,\underline b,
\overline a\,\overline b
\},
\max\{
\underline a\,\underline b,
\underline a\,\overline b,
\overline a\,\underline b,
\overline a\,\overline b
\}
\Bigr].
\]
\end{definition}

\begin{example}
If $A=[-1,3]$ and $B=[2,5]$, then the endpoint products are
\[
(-1)2=-2,\qquad (-1)5=-5,\qquad 3\cdot 2=6,\qquad 3\cdot 5=15.
\]
Hence
\[
A\cdot B=[-5,15].
\]
\end{example}

\begin{definition}[Interval division]
\cite{UnknownIntroductionIntervalAnalysis2002}
\label{def:interval-division}
For $A,B\in \mathcal{I}(\mathbb{R})$ with $0\notin B$ and
$B=[\underline b,\overline b]$, define
\[
A/B
=
\{\,a/b\mid a\in A,\ b\in B\,\}
=
A\cdot[\,1/\overline b,\ 1/\underline b\,].
\]
If $0\in B$, then $A/B$ is undefined in this elementary interval arithmetic setting.
\end{definition}

\begin{example}
If $A=[1,2]$ and $B=[4,5]$, then
\[
A/B
=
[1,2]\cdot [1/5,1/4]
=
[1/5,1/2].
\]
\end{example}
